\section{综合应用：让数学工具在一个真实问题中同时工作}

前面的学科分别讨论向量空间、概率、优化、数值误差、统计推断和信息量。真实建模却不会按课程边界出题。一个模型可能在纸面上有漂亮的最优解，却因矩阵病态而算不准；可能在验证集上预测良好，却没有可解释的因果含义；也可能优化完全收敛，却只是在精确拟合一个错误的损失函数。

综合应用的任务，就是把这些原本分开的判断重新放到同一条问题链上。

\subsection{一个模型必须连续回答的六个问题}

首先要明确希望改变或预测什么，数据又是怎样被观察到的。这个阶段决定样本单位、部署分布和评价目标，不能由算法代替。

接着要选择表示。线性代数告诉我们特征张成什么空间、哪些方向重复；降维与低秩建模进一步追问高维观测是否由少数潜在方向产生。

第三步是写出模型和损失。概率论区分联合、条件与边缘分布；线性模型展示平方误差、Bernoulli 似然和正则化如何从假设中产生；概率模型与隐变量则处理标签或状态无法直接观察的情形。

第四步是推断或优化。凸优化说明何时局部信息足以保证全局结论，矩阵微积分把复杂函数的梯度传回参数，采样与近似推断面对解析计算不可得的 posterior 和期望。

第五步是确认计算结果值得相信。数值分析区分问题本身的敏感性与算法误差；神经网络数学把 shape、反向传播、随机优化和浮点稳定性放在同一训练过程里。

最后才是统计评估。数理统计告诉我们估计量在重复抽样中怎样波动，数据切分与重抽样又在估计什么；信息论学习则从概率编码、最大熵、描述长度与表示压缩的角度审视目标。

这六步并非一次性流水线。残差结构可能迫使我们重新选择表示，数值不稳定可能暴露不可识别性，部署误差也可能说明训练分布或损失定义错了。建模是带反馈的推理过程。

\subsection{怎样选择阅读入口}

若希望建立一条完整的初学路线，从 线性模型 开始最合适。它的假设最透明，能同时看见几何、概率、优化和统计诊断。随后进入 降维，研究输入表示；再读 隐变量模型 与 近似推断，理解模型可写但 posterior 难算的情形。最后用 神经网络数学 和 信息论学习处理更灵活的表示与更抽象的学习目标。

也可以按问题进入：

\begin{itemize}
\tightlist
\item
  特征共线、模型系数不稳定：从线性模型的正则化开始；
\item
  高维数据需要压缩、去噪或补全：进入降维与低秩建模；
\item
  类别、簇或时间状态不可见：进入概率模型与隐变量；
\item
  积分、posterior 或抽样分布无法直接得到：进入采样与近似推断；
\item
  深度模型能运行却无法解释梯度或训练故障：进入神经网络数学；
\item
  想理解交叉熵、复杂度、压缩与表示之间的关系：进入信息论学习。
\end{itemize}

\subsection{贯穿全篇的判断边界}

在每个应用中，都应把下面几种``正确''分开：

\begin{itemize}
\tightlist
\item
  公式正确：推导在所写假设下成立；
\item
  数值正确：算法以足够精度求出了目标量；
\item
  统计正确：抽样、验证和不确定性分析与数据生成机制相符；
\item
  决策正确：指标与真实行动成本一致；
\item
  解释正确：预测相关性没有被夸大为因果或潜变量的唯一现实含义。
\end{itemize}

综合应用不是课程终点的算法合集，而是一次回到第一性问题的练习：我们究竟观察到了什么，额外假设了什么，计算了什么，又凭什么相信结论能用于尚未见到的世界。

